\documentclass[12pt]{article}
\usepackage[utf8]{inputenc}
\usepackage[T1]{fontenc}
\usepackage{hyperref}
\title{Bali Language Resources for K-12 Curriculum Development}
\author{Compiled by Claude Sonnet 5.5 (claude-sonnet-5-5) for Andr\'as Kornai}
\date{}
\begin{document}
\maketitle

\section{Introduction}
Bali (ISO 639-3 \texttt{bcn}; also Bibaali, Maya, Abaali, Ibaale) is an Atlantic-Congo (Central Adamawa, Yandangic) language of Nigeria, reported in Demsa, Adamawa State. Glottolog rates it as threatened/shifting. Speaker figures in sources are inconsistent (2,000 in one infobox, 100,000 in the article text). It must not be confused with Nga'ka (Bali Nyonga) of Cameroon, the Bali local government area of Taraba State, or the Fam language, which are different. Very little was found.

\section{Primary Educational Resources}
None found in search results.

\section{Linguistic and Reference Materials}
\begin{itemize}
\item Wikipedia, ``Bali language (Adamawa)'': \url{https://en.wikipedia.org/wiki/Bali_language_(Adamawa)}
\item Glottolog entry for ISO code bcn: \url{https://glottolog.org/resource/languoid/iso/bcn}
\item Glottolog classification page (bali1245): \url{https://glottolog.org/resource/languoid/id/bali1245.bigmap.html}
\end{itemize}

\section{Religious and Community Resources}
None found.

\section{Digital and Audio Resources}
Endangered Languages Project entry (listing alternate names; contents not verified): \url{https://endangeredlanguages.com/node/95062}.

\section{Curriculum Development Context}
No orthography, teaching materials or school use were found. Any curriculum work would begin with documentation and community consultation, and first confirm the speaker population and variety, given the conflicting figures and the many similarly named languages.
\end{document}